\section{The \texttt{MCViNE\_ConstantvQE} McStas Component}
Scatterer with a fixed momentum-transfer vector Q, weighted by a Gaussian in E.

\subsection*{Identification}
\begin{itemize}
  \item \textbf{Author:} Fahima Islam (McStas port of the MCViNE ConstantvQEKernel; MCViNE by J. Y. Y. Lin et al.)
  \item \textbf{Origin:} MCViNE, https://github.com/mcvine/mcvine (mccomponents/lib/kernels/sample/ConstantvQEKernel.cc)
  \item \textbf{Date:} 2026-09-29
\end{itemize}

\subsection*{Description}
kf = ki - Q (Q vector fixed in the sample frame); the event weight is multiplied by exp(-(E-E0)\textasciicircum{}2/(2 dE\textasciicircum{}2)) where E = Ei-Ef. Used for single-crystal resolution studies.

\textbf{Transport} (same as MCViNE HomogeneousNeutronScatterer): on the first passage the neutron is forced to scatter at a uniformly chosen depth, weighted by path length, attenuation exp(-(mu+sigma)x) and the scattering coefficient; on exit it is attenuated by exp(-(mu+sigma)L\_out). With order\textgreater{}1 further scatterings are sampled analogically (truncated exponential). With p\_transmit\textgreater{}0 that fraction of events is kept as the attenuated direct beam. mu(v) = absorption\_coefficient*2200/v. Events for which the kernel cannot scatter (kinematically forbidden) are absorbed.

Geometry: box (xwidth,yheight,zdepth), cylinder (radius,yheight), hollow cylinder (+thickness) or sphere (radius only). All vectors (Q, targets, reciprocal vectors, atom positions) are in the component's local frame.

Uses the shared runtime share/mcvine-lib.h/.c.

Example: MCViNE\_ConstantvQE(Qx=1, Qy=0, Qz=0.5, E=10, dE=1, scattering\_coefficient=10, xwidth=0.01, yheight=0.01, zdepth=0.01)

\subsection*{Input parameters}
Parameters in \textbf{boldface} are required; the others are optional.

\begin{longtable}{p{0.22\textwidth}p{0.12\textwidth}p{0.46\textwidth}p{0.14\textwidth}}
\toprule
\textbf{Name} & \textbf{Unit} & \textbf{Description} & \textbf{Default} \\
\midrule
\endhead
Qx & \AA{}$^{-1}$ & Q vector x & 0 \\
Qy & \AA{}$^{-1}$ & Q vector y & 0 \\
Qz & \AA{}$^{-1}$ & Q vector z & 0 \\
E & meV & Nominal energy transfer & 0 \\
dE & meV & Gaussian sigma of the energy weight & 1 \\
absorption\_coefficient & m$^{-1}$ & Absorption coefficient at 2200 m/s (MCViNE convention; scales as 1/v) & 0 \\
scattering\_coefficient & m$^{-1}$ & Scattering coefficient & 0 \\
sigma\_abs & barn & Alternative: absorption cross section per unit cell at 2200 m/s (used when Vc\textgreater{}0) & 0 \\
sigma\_scat & barn & Alternative: scattering cross section per unit cell (used when Vc\textgreater{}0) & 0 \\
Vc & \AA{}$^{3}$ & Unit cell volume; when \textgreater{}0, coefficients are computed from sigma\_abs, sigma\_scat & 0 \\
radius & m & Outer radius of a cylinder (with yheight) or of a sphere (yheight=0) & 0 \\
xwidth & m & Width of a box sample & 0 \\
yheight & m & Height of a box or cylinder sample & 0 \\
zdepth & m & Depth of a box sample & 0 \\
thickness & m & Wall thickness of a hollow cylinder (0: filled) & 0 \\
pack & 1 & Packing factor (scales absorption and scattering coefficients) & 1 \\
p\_transmit & 1 & Monte Carlo fraction of events kept as unscattered transmitted beam (0: always scatter) & 0 \\
order & 1 & Maximum number of scattering events per neutron (1: single scattering) & 1 \\
\bottomrule
\end{longtable}

\subsection*{Links}
\begin{itemize}
  \item Component source code found in file \texttt{MCViNE\_ConstantvQE.comp}.
  \item MCViNE documentation: https://mcvine.github.io
\end{itemize}
\IfFileExists{contrib/MCViNE_ConstantvQE_static.tex}{\input{contrib/MCViNE_ConstantvQE_static.tex}}{}